\documentclass[11pt, a4paper]{article}
\usepackage{graphicx}
\usepackage[T2A]{fontenc}
\usepackage[utf8]{inputenc}
\usepackage[russian]{babel}
\usepackage{amsmath, amssymb}
\usepackage[usenames]{color}
\usepackage{colortbl}
\usepackage{fancyhdr}
\usepackage{enumitem}
\usepackage{titlesec}
\usepackage{multicol}
\usepackage{tabularx}
\usepackage{caption}
\usepackage{tikz}
\usetikzlibrary{positioning, arrows.meta}
\usepackage[backend=biber, style=authoryear, maxcitenames=1, uniquelist=false]{biblatex}
\addbibresource{source.bib}
\usepackage{csquotes}
\usepackage[hidelinks]{hyperref}
\usepackage{xurl}
\usepackage{placeins}
\usepackage{longtable}
\usepackage{array}
\usepackage{pdflscape}

\usepackage{geometry}
\geometry{
  a4paper,
  top=30mm,
  right=5mm,
  bottom=20mm,
  left=5mm,
  headheight=61pt
}
\setlength{\columnsep}{8mm}
\raggedcolumns  % не растягивать колонки по высоте
\emergencystretch=2em
\renewcommand{\topfraction}{0.95}
\renewcommand{\bottomfraction}{0.6}
\renewcommand{\textfraction}{0.05}
\renewcommand{\floatpagefraction}{0.85}
\setcounter{topnumber}{3}\setcounter{totalnumber}{4}

\newcommand{\mydate}{Октябрь, 2026}
\newcommand{\myauthor}{В. Алексеев, С. Зинина, И. Коновалов, В. Горбатенко, А. Соколова}

\fancyhf{}
\fancyhead[L]{\large \textit{Научная студия 2026. Вычислительная онкология. Общий анализ данных.}}
\fancyfoot[C]{\thepage}
\pagestyle{fancy}

\fancypagestyle{firstpage}{
  \fancyhf{}
  \fancyhead[L]{\hspace{-15pt}\includegraphics[width=0.3\linewidth]{fig/cu_logo.png}}
  \fancyhead[R]{\mydate\\\textit{\myauthor}\\}
  \fancyfoot[C]{\thepage}
}

\newcommand{\code}[1]{\texttt{\small #1}}

\begin{document}
\thispagestyle{firstpage}

\begin{multicols}{2}
\section*{Материалы и методы}

\subsection*{Исход и момент применения модели}

Модель предсказывает локальный рецидив первичной опухоли после радикальной лучевой терапии (ЛТ). Целевая переменная \code{Target\_local\_recurrence} равна 1, если за период наблюдения зарегистрирован рецидив в той же или иной подобласти ротоглотки, и 0, если рецидива не было; значение перенесено без изменений из \code{Local\_tumor\_recurrence}, заполненной авторами набора по данным наблюдения \autocite{Elhalawani2017}. Эта конечная точка выбрана потому, что локальный контроль -- прямая цель облучения первичного очага.\vspace{6pt}\\
В таблицах нет дат, длительности наблюдения и времени до рецидива, поэтому исход нельзя пересчитать на иной горизонт и он используется в форме, заданной организаторами. Горизонт задан критерием включения -- наблюдением более двух лет \autocite{Elhalawani2017}, покрывающим период, на который приходится основная часть локальных неудач после ЛТ. По критериям отбора все 128 пациентов без события наблюдались не меньше горизонта, однако индивидуально проверить это по таблице нельзя.\vspace{6pt}\\
У 158 пациентов тестовой части исход скрыт организаторами, т.\,е. неизвестен, а не отрицателен; присвоение им нуля занизило бы частоту рецидива, поэтому они исключены целиком. В анализ вошли 140 пациентов обучающей части: 12 с рецидивом и 128 без него. Малое число событий -- главное ограничение данных, требующее жёсткого ограничения числа признаков при моделировании \autocite{Riley2020}.\vspace{6pt}\\
Модель применяется после стадирования и утверждения плана лечения, до начала любой его части. К этому моменту известны диагностическая КТ, по которой построены маски GTVp и GTVn, демографические данные, анамнез курения, TNM и стадия, степень дифференцировки по биопсии, предписанные доза и фракционирование и решение о химиотерапии. Оценка риска при выбранном режиме позволяет обсуждать эскалацию дозы или более плотное наблюдение после ЛТ; раньше параметры лечения неизвестны, позже прогноз теряет практическую ценность. Поэтому как признаки нельзя использовать фактическую длительность курса \code{Radiation\_treatment\_course\_duration}, зависящую от перерывов и известную лишь после окончания ЛТ, и жизненный статус \code{KM\_Overall\_survival\_censor}, относящийся к периоду после лечения: это была бы утечка информации из будущего \autocite{Kaufman2012}.

\subsection*{Первичный анализ данных}

Использована итоговая таблица этапа извлечения признаков (298 пациентов, 235 столбцов: 20 клинических и служебных, число узлов \code{GTVn\_n\_nodes} и по 107 признаков GTVp и GTVn); исходный файл не изменялся. Для 14 категориальных переменных выведены все значения, для количественных -- наименьшее и наибольшее, для клинических -- в сравнении с допустимым диапазоном. Разных написаний одной категории и значений вне диапазонов нет, пропуски записаны пустыми ячейками, а в \code{N\_category} и \code{AJCC\_Stage} встречается служебный код X (<<не может быть оценена>>). Дат в наборе нет, поэтому порядок событий и длительности не проверялись.\vspace{6pt}\\
Значение заменялось верным, только если оно однозначно восстанавливается по другим полям записи; служебные коды и невосстановимые противоречия заменялись пропуском. Код X у одного пациента заменён пропуском в обеих переменных. У одного пациента при T4N1 указана стадия III, хотя по 7-му изданию AJCC при M0 категория T4 означает стадию IV \autocite{Edge2010}; стадия выводится из T и N однозначно и исправлена. У 7 пациентов со статусом Current или Former указано 0 пачко-лет, что противоречит статусу; экспозицию восстановить нельзя, значения заменены пропуском. Всего изменено 10 значений: 2 в \code{AJCC\_Stage}, 1 в \code{N\_category} и 7 в \code{Smoking\_Pack-Years}.\vspace{6pt}\\
Доза на фракцию (1,8--2,4~Гр) и число суток на фракцию (0,98--1,70; меньше единицы при облучении дважды в день) правдоподобны. Плотность ниже $-1024$~HU у 12 пациентов и выше 3071~HU у 5 объясняется выбросами B-сплайн-интерполяции и артефактами от металла; это свойства вычисленных признаков, а не ошибки записи, и они сохранены. Маска GTVn есть у 2 пациентов с N0 и отсутствует у 19 с N1--N3; это не исправлялось, так как маски строились по видимым на КТ узлам и отсутствуют после эксцизионной биопсии узла \autocite{Elhalawani2017}.\vspace{6pt}\\
Пропуски есть в \code{Local\_tumor\_recurrence} (158), \code{Pathological\_grade} (43), \code{Smoking\_Pack-Years} (21, из них 7 после исправления), \code{Race}, \code{N\_category} и \code{AJCC\_Stage} (по 1), а также во всех 107 признаках GTVn у 45 пациентов; в GTVp пропусков нет. Пропуски GTVn приходятся ровно на пациентов с \code{GTVn\_n\_nodes}~=~0, т.\,е. означают отсутствие узлов, а не потерю данных. В когорте с известным исходом доля пропусков составляет 11,4~\% в степени дифференцировки, 7,1~\% в пачко-годах и 14,3~\% в GTVn. Переменная исключалась при доле пропусков выше 30~\%, так как иначе её значения определялись бы в основном моделью заполнения; таких переменных нет.\vspace{6pt}\\
Исключены идентификатор (сохранён как ключ), источник целевой переменной, служебная \code{Fold}, две недоступные в момент применения переменные и 2 дубликата: \code{firstorder\_TotalEnergy} в обеих масках совпадает с \code{firstorder\_Energy}, поскольку при вокселе 1~мм$^3$ они отличаются лишь множителем \autocite{PyRadiomicsDocs}. Констант нет. Добавлена вычисленная \code{Dose\_per\_fraction}, так как при равной суммарной дозе эффект облучения зависит от фракционирования.\vspace{6pt}\\
Избыточность 218 количественных признаков сокращалась последовательным отбором: признак оставлялся, если модуль корреляции с каждым ранее оставленным не превышал 0,7, иначе исключался и засчитывался наиболее связанному оставленному. Порядок предпочтения не зависит от исхода: клинические признаки как доступные в рутинной практике, затем GTVp как область конечной точки, затем число узлов и GTVn; внутри маски первыми идут объём (установленный прогностический фактор) и сферичность, далее порядок PyRadiomics. Использовалась ранговая корреляция Спирмена, так как распределения признаков сильно скошены и содержат выбросы плотности. Порог 0,7 (более половины общей ранговой дисперсии) выбран жёстким из-за малого числа событий. Корреляции считались по всем 298 пациентам без использования исхода, для GTVn -- по 253 пациентам с узлами. Оставлен 41 признак (4 клинических, \code{GTVn\_n\_nodes}, 16 GTVp и 20 GTVn), исключено 177, наибольший модуль корреляции среди оставленных -- 0,698. Больше всего признаков заменяют энтропия GTVn (32) и GTVp (26), минимальная плотность GTVp (25) и объём GTVp (23).\vspace{6pt}\\
Аналитическая таблица содержит 140 пациентов, код пациента, целевую переменную и 52 признака: 11 категориальных и 41 количественный. Пропуски (16 в степени дифференцировки, 10 в пачко-годах, по 20 в признаках GTVn) на этом этапе не заполнялись, а связь признаков с исходом не изучалась, чтобы заполнение и отбор по исходу выполнять внутри кросс-валидации без оптимистичного смещения оценки.

\end{multicols}

\FloatBarrier
\begin{multicols}{2}
\subsection*{Ссылки}
\printbibliography[heading=none]
\end{multicols}

\newpage
\begin{landscape}
\subsection*{Словарь переменных и решений}

\noindent Приведены все клинические и служебные переменные исходной таблицы, число узлов \code{GTVn\_n\_nodes}, вычисленная \code{Dose\_per\_fraction}, целевая переменная и оставленные радиомические признаки (16 из 107 для GTVp и 20 из 107 для GTVn) с числом исключённых признаков, которые каждый из них заменяет. Пропуски указаны после исправления значений во всей таблице (298 пациентов) и в аналитической когорте (140 пациентов); момент доступности -- относительно применения модели после утверждения плана лечения.

\footnotesize
\setlength{\tabcolsep}{3.5pt}
\renewcommand{\arraystretch}{1.15}
\begin{longtable}{>{\raggedright\arraybackslash}p{5.0cm} >{\raggedright\arraybackslash}p{4.4cm} >{\raggedright\arraybackslash}p{3.6cm} >{\raggedright\arraybackslash}p{3.1cm} >{\raggedright\arraybackslash}p{2.2cm} >{\raggedright\arraybackslash}p{4.9cm}}
\hline\hline
Переменная & Смысл & Тип, допустимые значения & Момент доступности & Пропуски (298 / 140) & Решение и комментарий \\ \hline
\endfirsthead
\hline\hline
Переменная & Смысл & Тип, допустимые значения & Момент доступности & Пропуски (298 / 140) & Решение и комментарий \\ \hline
\endhead
\hline
\multicolumn{6}{r}{\textit{Продолжение на следующей странице}}\\
\endfoot
\hline\hline
\caption{Словарь переменных и решений первичного анализа.}
\endlastfoot
\texttt{Patient\_\allowbreak{}ID} & код пациента, присвоенный при анонимизации, совпадает с именем папки снимков & целый, 1–315, уникален & служебная, всегда & 0 / 0 & оставить как ключ строки, не признак: идентификатор не несёт клинической информации \\
\texttt{Fold} & часть набора соревнования & категориальный: Training, Test & служебная, задана организаторами & 0 / 0 & исключить: служебная переменная \\
\texttt{Local\_\allowbreak{}tumor\_\allowbreak{}recurrence} & локальный рецидив первичной опухоли за период наблюдения & бинарный: 0, 1 & после наблюдения не менее 2 лет & 158 (53,0~\%) / 0 & исключить из признаков: источник целевой переменной \\
\texttt{Gender} & пол & категориальный: Male, Female & до начала лечения, при постановке диагноза & 0 / 0 & оставить \\
\texttt{Age\_\allowbreak{}at\_\allowbreak{}diagnosis} & возраст на момент диагноза, лет & целый, 30–89 & до начала лечения, при постановке диагноза & 0 / 0 & оставить, заменяет 0 исключённых признаков \\
\texttt{Race} & раса & категориальный: White, Black, Hispanic, Asian, American\_Indian/\allowbreak{}Alaska\_Native & до начала лечения, при постановке диагноза & 1 (0,3~\%) / 0 & оставить, пустая ячейка~--- пропуск, не заполнять \\
\texttt{Tumor\_\allowbreak{}side} & латеральность опухоли & категориальный: L, R, Bilateral & до начала лечения, после стадирования & 0 / 0 & оставить \\
\texttt{Tumor\_\allowbreak{}subsite} & подобласть ротоглотки & категориальный: BOT, Tonsil, GPS, Soft\_palate, Pharyngeal\_wall, Other & до начала лечения, после стадирования & 0 / 0 & оставить \\
\texttt{T\_\allowbreak{}category} & категория T (AJCC 7) & порядковый: 1–4 & до начала лечения, после стадирования & 0 / 0 & оставить \\
\texttt{N\_\allowbreak{}category} & категория N (AJCC 7) & порядковый: 0, 1, 2a, 2b, 2c, 3 & до начала лечения, после стадирования & 1 (0,3~\%) / 0 & оставить, служебный код X (1 значение) заменён пропуском \\
\texttt{AJCC\_\allowbreak{}Stage} & стадия по AJCC 7, определяется T и N при M0 & порядковый: I–IV & до начала лечения, после стадирования & 1 (0,3~\%) / 0 & оставить, X (1) заменён пропуском, III при T4N1 (1) исправлено на IV \\
\texttt{Pathological\_\allowbreak{}grade} & степень дифференцировки опухоли & порядковый: I, I-II, II, II-III, III, IV & до начала лечения, по биопсии & 43 (14,4~\%) / 16 (11,4~\%) & оставить, пропуски не заполнять, доля ниже порога 30~\% \\
\texttt{Smoking\_\allowbreak{}status\_\allowbreak{}at\_\allowbreak{}diagnosis} & статус курения на момент диагноза & категориальный: Never, Former, Current & до начала лечения, при постановке диагноза & 0 / 0 & оставить \\
\texttt{Smoking\_\allowbreak{}Pack-Years} & кумулятивное курение, пачко-лет & непрерывный, $\geq 0$, 0 только при Never & до начала лечения, при постановке диагноза & 21 (7,0~\%) / 10 (7,1~\%) & оставить, 7 нулей при Current/Former заменены пропуском, заменяет 0 исключённых признаков \\
\texttt{Radiation\_\allowbreak{}treatment\_\allowbreak{}course\_\allowbreak{}duration} & фактическая длительность курса ЛТ, дни & целый, 32–56 & после окончания ЛТ & 0 / 0 & исключить: неизвестна в момент применения модели \\
\texttt{Total\_\allowbreak{}prescribed\_\allowbreak{}Radiation\_\allowbreak{}treatment\_\allowbreak{}dose} & суммарная предписанная доза, Гр & целый: 60, 63, 66, 70, 72 & до начала ЛТ, при утверждении плана лечения & 0 / 0 & оставить, заменяет 0 исключённых признаков \\
\texttt{\#\_\allowbreak{}Radiation\_\allowbreak{}treatment\_\allowbreak{}fractions} & число фракций & целый, 30–40 & до начала ЛТ, при утверждении плана лечения & 0 / 0 & исключить: избыточна, $|\rho| > 0{,}7$ с Dose\_per\_fraction \\
\texttt{Induction\_\allowbreak{}Chemotherapy} & индукционная химиотерапия до ЛТ & бинарный: Y, N & до начала ЛТ, при утверждении плана лечения & 0 / 0 & оставить \\
\texttt{Concurrent\_\allowbreak{}chemotherapy} & химиотерапия одновременно с ЛТ & бинарный: Y, N & до начала ЛТ, при утверждении плана лечения & 0 / 0 & оставить \\
\texttt{KM\_\allowbreak{}Overall\_\allowbreak{}survival\_\allowbreak{}censor} & жизненный статус на последнем наблюдении & бинарный: 1~--- жив, 0~--- умер & после наблюдения & 0 / 0 & исключить: альтернативный исход, известен в обеих частях набора, но без времени наблюдения, неизвестен в момент применения модели \\
\texttt{GTVn\_\allowbreak{}n\_\allowbreak{}nodes} & число поражённых лимфоузлов, выделенных на маске GTVn & целый, 0–16 & до начала лечения, по диагностической КТ & 0 / 0 & исключить вместе со всеми признаками GTVn: по условиям соревнования модель строится по первичной опухоли \\
\texttt{Dose\_\allowbreak{}per\_\allowbreak{}fraction} & вычисленная: доза на фракцию, Гр (доза / число фракций) & непрерывный, 1,8–2,4 & до начала ЛТ, при утверждении плана лечения & 0 / 0 & оставить, заменяет 1 исключённый признак \\
\texttt{Target\_\allowbreak{}local\_\allowbreak{}recurrence} & целевая: 1~--- локальный рецидив, 0~--- нет & бинарный: 0 (128), 1 (12) & после наблюдения не менее 2 лет & 158 (53,0~\%) / 0 & целевая переменная \\
\texttt{GTVp\_\allowbreak{}shape\_\allowbreak{}Mesh\allowbreak{}Volume} & объём по треугольной сетке (GTVp) & непрерывный, мм$^3$, $>0$ & до начала лечения, по диагностической КТ & 0 / 0 & оставить, заменяет 23 исключённых признака \\
\texttt{GTVp\_\allowbreak{}shape\_\allowbreak{}Sphericity} & сферичность (GTVp) & непрерывный, (0, 1] & до начала лечения, по диагностической КТ & 0 / 0 & оставить, заменяет 0 исключённых признаков \\
\texttt{GTVp\_\allowbreak{}shape\_\allowbreak{}Elongation} & вытянутость: отношение второй главной оси к первой (GTVp) & непрерывный, (0, 1] & до начала лечения, по диагностической КТ & 0 / 0 & оставить, заменяет 0 исключённых признаков \\
\texttt{GTVp\_\allowbreak{}shape\_\allowbreak{}Flatness} & уплощённость: отношение третьей главной оси к первой (GTVp) & непрерывный, (0, 1] & до начала лечения, по диагностической КТ & 0 / 0 & оставить, заменяет 0 исключённых признаков \\
\texttt{GTVp\_\allowbreak{}firstorder\_\allowbreak{}10Percentile} & 10-й перцентиль плотности (GTVp) & непрерывный, HU & до начала лечения, по диагностической КТ & 0 / 0 & оставить, заменяет 1 исключённый признак \\
\texttt{GTVp\_\allowbreak{}firstorder\_\allowbreak{}90Percentile} & 90-й перцентиль плотности (GTVp) & непрерывный, HU & до начала лечения, по диагностической КТ & 0 / 0 & оставить, заменяет 2 исключённых признака \\
\texttt{GTVp\_\allowbreak{}firstorder\_\allowbreak{}Entropy} & энтропия гистограммы плотностей (GTVp) & непрерывный, бит, $\geq 0$ & до начала лечения, по диагностической КТ & 0 / 0 & оставить, заменяет 26 исключённых признаков \\
\texttt{GTVp\_\allowbreak{}firstorder\_\allowbreak{}Kurtosis} & эксцесс распределения плотностей (GTVp) & непрерывный, $\geq 1$ & до начала лечения, по диагностической КТ & 0 / 0 & оставить, заменяет 3 исключённых признака \\
\texttt{GTVp\_\allowbreak{}firstorder\_\allowbreak{}Maximum} & максимальная плотность (GTVp) & непрерывный, HU & до начала лечения, по диагностической КТ & 0 / 0 & оставить, заменяет 0 исключённых признаков \\
\texttt{GTVp\_\allowbreak{}firstorder\_\allowbreak{}Minimum} & минимальная плотность (GTVp) & непрерывный, HU & до начала лечения, по диагностической КТ & 0 / 0 & оставить, заменяет 25 исключённых признаков \\
\texttt{GTVp\_\allowbreak{}glcm\_\allowbreak{}Correlation} & GLCM: корреляция уровней соседних вокселей (GTVp) & непрерывный, $[-1, 1]$ & до начала лечения, по диагностической КТ & 0 / 0 & оставить, заменяет 4 исключённых признака \\
\texttt{GTVp\_\allowbreak{}glcm\_\allowbreak{}Inverse\allowbreak{}Variance} & GLCM: обратная дисперсия (GTVp) & непрерывный, $\geq 0$ & до начала лечения, по диагностической КТ & 0 / 0 & оставить, заменяет 0 исключённых признаков \\
\texttt{GTVp\_\allowbreak{}gldm\_\allowbreak{}Dependence\allowbreak{}Non\allowbreak{}Uniformity\allowbreak{}Normalized} & GLDM: нормированная неоднородность зависимостей (GTVp) & непрерывный, (0, 1] & до начала лечения, по диагностической КТ & 0 / 0 & оставить, заменяет 5 исключённых признаков \\
\texttt{GTVp\_\allowbreak{}glszm\_\allowbreak{}Size\allowbreak{}Zone\allowbreak{}Non\allowbreak{}Uniformity\allowbreak{}Normalized} & GLSZM: нормированная неоднородность размеров зон (GTVp) & непрерывный, (0, 1] & до начала лечения, по диагностической КТ & 0 / 0 & оставить, заменяет 1 исключённый признак \\
\texttt{GTVp\_\allowbreak{}ngtdm\_\allowbreak{}Busyness} & NGTDM: «занятость»~--- частота смены тонов (GTVp) & непрерывный, $\geq 0$ & до начала лечения, по диагностической КТ & 0 / 0 & оставить, заменяет 1 исключённый признак \\
\texttt{GTVp\_\allowbreak{}ngtdm\_\allowbreak{}Contrast} & NGTDM: контраст соседних тонов (GTVp) & непрерывный, $\geq 0$ & до начала лечения, по диагностической КТ & 0 / 0 & оставить, заменяет 0 исключённых признаков \\

\end{longtable}
\end{landscape}

\end{document}
